\documentclass{article}

\usepackage[top=0.5in, bottom=0.5in, left=0.5in, right=0.5in]{geometry}
\usepackage{enumitem}
\usepackage[utf8]{inputenc}
\usepackage[english]{babel}
\usepackage{multicol}
\usepackage{fontawesome}
\usepackage{hyperref}

\setlength{\parindent}{0pt}



\begin{document}
\begin{center}
\thispagestyle{empty}
\begin{Large}
    \textbf{Robin Uhrich\\}
\end{Large}
\begin{center}   
    \begin{tabular}{c l l c l}
         \normalsize \faMapMarker & University of Tübingen & \hfill & \faGraduationCap &  M.Sc. Machine Learning \\ 
         \normalsize \faAt & \href{mailto:robin.uhrich@gmail.com}{robin.uhrich@gmail.com}
         & \hfil & \normalsize \faPhone & +49 172 9606083 \\
         \normalsize \faGlobe & \href{https://robinu434.github.io/}{robinu434.github.io} & \hfil & & \\
         \normalsize \faGithubSquare &  \href{https://github.com/RobinU434/}{RobinU434} & \hfil & \normalsize \faLinkedinSquare & \href{https://www.linkedin.com/in/robin-uhrich-bb2891293}{robin-uhrich-bb2891293}
    \end{tabular}
\end{center}
    \hrulefill
\end{center}


%%%%%%%%%%%%%%%%%%%%%%%%%%%%%%%%%%%%%%%%%%%%%%%%%%%%%%%%%%%%%%%
% OBJECTIVE
% Who you are, what domain, what are you looking for and when?
%%%%%%%%%%%%%%%%%%%%%%%%%%%%%%%%%%%%%%%%%%%%%%%%%%%%%%%%%%%%%%%
%\noindent \textbf{\underline{OBJECTIVE}} \\
%\noindent I am an ungraduated student in computer science in the second year at the Albert Ludwigs Universität Freiburg. I \\
% Skill 1 (level of expertise), Skill 2 (level of expertise), Skill 3 (level of expertise) \\

%%%%%%%%%%%%%%%%%%%%%%%%%%%%%%%%%%%%%%%%%%%%%%%%%%%%%%%%%%%%%%%
% EDUCATION
% University name, degree, year of graduation, GPA (optional)
%%%%%%%%%%%%%%%%%%%%%%%%%%%%%%%%%%%%%%%%%%%%%%%%%%%%%%%%%%%%%%%
\noindent \textbf{\underline{EDUCATION}} \\

% \noindent\textbf{High School Backnang} \hfill Backnang, Germany \\
% \textit{Technical High School} \hfill September 2015 $-$ July 2018 \\ \\
\noindent\textbf{University of Tübingen} \hfill Tübingen, Germany  \\
\textit{Master Machine Learning} \hfill since October 2023 \\ \\
\textbf{Politecnico di Milano} \hfill Milano, Italy  \\
\textit{Erasmus Exchange, Master Computer Science} \hfill February 2025 $-$ August 2025 \\ \\
\textbf{University of Freiburg}\hfill Freiburg im Breisgau, Germany  \\
\textit{Bachelor Computer Science} \hfill October 2019 $-$ September 2023 \\ \\
\textbf{University of Stuttgart} \hfill Stuttgart, Germany \\
\textit{Aeronautics and Space Engineering} \hfill October 2018 $-$ September 2019 \\  \\


%%%%%%%%%%%%%%%%%%%%%%%%%%%%%%%%%%%%%%%%%%%%%%%%%%%%%%%%%%%%%%%
% PUBLICATIONS
%%%%%%%%%%%%%%%%%%%%%%%%%%%%%%%%%%%%%%%%%%%%%%%%%%%%%%%%%%%%%%%
\noindent \textbf{\underline{PUBLICATIONS}} \\

\noindent \textbf{Occupancy Modeling via Drifting: One-Step Generation Enables Direct Policy Extraction} \hfill 2026 \\
\underline{Robin Uhrich}, Marco Bagatella, Núria Armengol Urpí, Pavel Kolev, Georg Martius \\
\textit{Preprint}, \href{https://arxiv.org/abs/2610.19883}{arXiv:2610.19883}
\vspace{0.4cm}\\

%%%%%%%%%%%%%%%%%%%%%%%%%%%%%%%%%%%%%%%%%%%%%%%%%%%%%%%%%%%%%%%
% PROJECT
% What did you do?
% How did you do it? -> Skills and technologies
% What impact did you create? -> Numbers and percentages.
%
% Talk about team work, initiative, soft skills.
%
% Can also include personal projects, competitions, contribution to Open source.
%%%%%%%%%%%%%%%%%%%%%%%%%%%%%%%%%%%%%%%%%%%%%%%%%%%%%%%%%%%%%%%
\noindent \textbf{\underline{PROJECT WORK}} \\

% \noindent \textbf{Umbrella Satellite Antenna} \textit{High School Project} \hfill  October 2016 $-$ July 2017
% \begin{itemize}[noitemsep,nolistsep,leftmargin=*]
%     \item {Collaboration with TESAT Spacecom GmbH.}
%     \item {developing a high compact parabolic satellite antenna using an umbrella}
%     \item {using rapid prototyping and mechanical CAD systems\\}
% \end{itemize}


\textbf{University} \\


% \noindent \textbf{Wall$\land$E} \textit{System Design Project} \hfill  September 2019 $-$ February 2020
% \begin{itemize}[noitemsep,nolistsep,leftmargin=*]
%     \item {Course project on developing a robot using Lego EV3 set \\ 
%     to solve multi-task problem settings.}
%     \item {Programming directly on the EV3 computer in $\mu$-Python. }
%     \item Implementation of a PID controller for line following based \\ on light and ultrasonic sensors.
%     \item Design and implementation of a finite state machine for different tasks.\\
% \end{itemize}
% 
% \noindent \textbf{2048} \textit{Game development in C++ class} \hfill  July 2020 $-$ August 2020
% \begin{itemize}[noitemsep,nolistsep,leftmargin=*]
%     \item Course project within the lecture on "Advanced Programming" \\ 
%     organized by Prof. Hannah Bast
%     \item Implemented an object-oriented version of the game "2048" in C++.
%     \item Implementation of a CI/CD pipeline, including tests and documentation.\\
% \end{itemize}

\noindent \textbf{Master Thesis} \hfill since October 2025
\begin{itemize}[noitemsep,nolistsep,leftmargin=*]
    \item Master thesis at the lab of Prof. Georg Martius, supervised by \\
    Dr. Pavel Kolev, Núria Armengol Urpí and Marco Bagatella.
    \item Topics: occupancy modelling, generative models, and offline meta-RL.
    \item Resulted in the preprint \href{https://arxiv.org/abs/2610.19883}{arXiv:2610.19883} (see Publications).\\
\end{itemize}

\noindent \textbf{Distributed RL on Optimal Transport} \textit{Research Project} \href{https://github.com/RobinU434/SwitchFL}{\faGithubSquare} \hfill April 2025 $-$ October 2025
\begin{itemize}[noitemsep,nolistsep,leftmargin=*]
    \item Project done at the lab of Prof. Marcello Restelli, \\
    supervised by Dr. Marco Mussi.
    \item Redesign of the rail environment from Flatland3 for multi-agent \\
    RL in an asynchronous setting to solve the train-scheduling problem.
    \item Switched the perspective from train-focused to switch-focused.
    \item Tools used: Python, Flatland3, Gymnasium, PettingZoo.\\
\end{itemize}

\noindent \textbf{Field Robot Event} \textit{Multi-Disciplinary Project at PoliMi}\hfill March 2025 $-$ June 2025
\begin{itemize}[noitemsep,nolistsep,leftmargin=*]
    \item Course project on robotics under Prof. Matteo Matteucci, \\
    Dr. Simone Mentasti and Mirko Usuelli.
    \item Third place at the Field Robot Event 2025 competition.
    \item Building and programming a robot for an \\
    agricultural-robotics competition in a team of nine.
    \item Implementation of navigation, computer-vision \\ 
    and decision making algorithms.
    \item Conducted study on visual object tracking in a \\
    moving stereo-camera setting with DeepSORT. 
    \item Tools used: Python, ROS2, PyTorch, Yolo\\
\end{itemize}

\noindent \textbf{Multi-Agent RL in Laser Hockey} \textit{Course Project} \href{https://github.com/RobinU434/LaserHockeyProject}{\faGithubSquare} \hfill December 2024 $-$ February 2025
\begin{itemize}[noitemsep,nolistsep,leftmargin=*]
    \item Course project on ``Reinforcement Learning'' taught by Prof. Georg Martius.
    \item Took part in an organized multi-agent RL competition in laser hockey.
    \item Extended SAC from Stable-Baselines3 with a custom replay buffer.
    \item Implemented and extended Deep Dyna-Q with entropy regularization.
    \item Contributed to the development of the laser-hockey environment.\\
\end{itemize}

\noindent \textbf{Generalized Simulation Based Inference} \textit{Practical Machine Learning} 
\href{https://github.com/mackelab/neuralgbi_diffusion}{\faGithubSquare} \hfill October 2024 $-$ April 2025
\begin{itemize}[noitemsep,nolistsep,leftmargin=*]
    \item ML-Lab project under supervision of Prof. Jakob Macke and Julius Vetter.
    \item Extended sampling strategies in simulation-based inference with \\
    amortized cost estimation presented by Dr. Richard Gao (\href{https://arxiv.org/abs/2305.15208}{paper}) \\
    with guided diffusion models.
    \item Re-implemented GBI-ACE and used it for benchmarking \\ 
    on toy and real world datasets.
    \item Implementation was done using the PyTorch framework for deep learning.\\
\end{itemize}

\noindent \textbf{LR - Scheduler Analysis} \textit{Seminar Project} \href{https://github.com/RobinU434/LRSchedulerAnalysis}{\faGithubSquare}\hfill April 2024 $-$ July 2024
\begin{itemize}[noitemsep,nolistsep,leftmargin=*]
    \item Project work for the seminar: "Deep Learning Research Kitchen" \\ 
    hosted by Prof. Philipp Hennig, Dr. Frank Schneider, \\
    Dr. Jonas Geiping, Dr. Antonio Orvieto.
    \item Conducted a study of multiple learning-rate schedulers and how their \\
    hyperparameters affect the training of a Vision Transformer \\
    trained on CIFAR-10.
    \item Implemented pipeline using PyTorch-Lightning, TensorBoard, \\
    Numpy and Matplotlib for training and analysis. \\ 
\end{itemize}


\noindent \textbf{Precipitation Prediction Analysis} \textit{Data Literacy Group Project} \href{https://github.com/RobinU434/PrecipitationPredictionAnalysis}{\faGithubSquare}\hfill December 2023 $-$ February 2024
\begin{itemize}[noitemsep,nolistsep,leftmargin=*]
    \item Course project on "Data Literacy" by Prof. Philipp Hennig. \\
    Project was supervised by Emilia Magnani. 
    \item Data science project on climate data in a group of four.
    \item Performed steps: problem identification, \\ data collection, and data analysis.
    \item Analyzed the accuracy of the precipitation forecast \\
    of the German weather agency against real data for the \\
    state of Baden-Württemberg, Germany.\\
\end{itemize}

\noindent \textbf{EASRL-LM-2D-IKB} \textit{Bachelor Thesis} \href{https://github.com/RobinU434/IK-RL}{\faGithubSquare \ environment} \href{https://github.com/RobinU434/SPARK/tree/v0.0.1-bachelor-thesis-submission}{\faGithubSquare \ pipeline}\hfill July 2023 $-$ September 2023
\begin{itemize}[noitemsep,nolistsep,leftmargin=*]
    \item Bachelor thesis at the lab of Prof. Joschka Boedecker, \\ 
    supervised by Jasper Hoffmann.  
    \item Topic: Expanding Action Space in Reinforcement Learning through \\ Latent Models: A 2D Inverse Kinematics Benchmark.
    \item Merging Soft Actor-Critic with generative models such as (pretrained) \\
    Variational Autoencoder for noise augmentation and faster learning.
    \item Implementation of a custom inverse kinematics environment \\ 
    with a scalable action space. \\
\end{itemize} 


\noindent \textbf{Learning Quadrupedal Locomotion From Scratch} \textit{Bachelor Project} \href{https://github.com/feinran/locomotion_simulation}{\faGithubSquare}\hfill  April 2022 $-$ September 2022
\begin{itemize}[noitemsep,nolistsep,leftmargin=*]
    \item Bachelor research project together with Felix Schulz \\ at the lab of Prof. Joschka Boedecker. Supervised by Jasper Hoffmann.
    \item Implemented training pipeline for reinforcement learning \\ 
    experiments around PPO concerned with perturbation and reward-shaping.
    \item Minimalist re-implementation of \href{https://arxiv.org/abs/2109.14549}{\textit{Vision-Guided Quadrupedal Locomotion}}.
    \item Used and extended an already existing physics-simulator \\ 
    which is based on PyBullet.
    \item Main libraries: Stable-Baselines3, PyTorch, Gym, \\ Weights \& Biases, and TensorBoard.
    \item Conducted a study on closing the sim-to-real gap on the Unitree A1 robot. \\
\end{itemize}


\noindent \textbf{ALE Plots} \textit{Seminar Project} \href{https://github.com/RobinU434/AccumulatedLocalEffectPlots}{\faGithubSquare} \hfill  May 2022 $-$ September 2022
\begin{itemize}[noitemsep,nolistsep,leftmargin=*]
    \item Project work for the seminar: "Fair and Interpretable Machine Learning" \\ 
    hosted by Dr. Janek Thomas.
    \item Re-implementation of \href{https://academic.oup.com/jrsssb/article/82/4/1059/7056085}{accumulated local effect plots}.
    \item Implementation was done using Python alongside Numpy and Matplotlib. \\
\end{itemize}


\noindent \textbf{Humboldt} \textit{Game development in class ``software practical training''} \hfill  October 2020 $-$ April 2021
\begin{itemize}[noitemsep,nolistsep,leftmargin=*]
    \item Course project on ``Practical Software Development'' \\
    organized by Prof. Andreas Podelski and Dr. Vincent Langenfeld. \\
    Supervised by Antonia Hinkerohe.
    \item Designed, implemented, and tested the real-time \\
    multi-agent strategy game ``Humboldt''.
    \item Collaborated in a team of six and used Git for \\ 
    version control.
    \item Programmed in C\# using the XNA-(monogame)-framework. \\
\end{itemize}



\vspace{0.5cm}
\textbf{Student Assistant} \\

\noindent \textbf{Human2Humanoid VLA} \textit{Research Project} \hfill October 2025 $-$ March 2026
\begin{itemize}[noitemsep,nolistsep,leftmargin=*]
    \item Project done at the lab of Prof. Georg Martius and Prof. Wieland Brendel, \\
    under the supervision of Anna Manasyan and Jens Beißwenger.
    \item Finetuning Vision-Language-Action models applied to humanoid robotics.
    \item Collecting humanoid demonstration data. 
    \item Conducting studies on action representations and generalization in behavior cloning.\\
\end{itemize}

\noindent \textbf{Hand Kinematics Estimation from EMG Data} \textit{Software Development} \hfill August 2024 $-$ February 2025
\begin{itemize}[noitemsep,nolistsep,leftmargin=*]
    \item Project done at BG-Clinic under the supervision of \\ 
    Daniel Höglinger at the lab of Dr. Cosima Prahm.
    \item Designed a deep learning pipeline for \\ hand kinematics estimation based on EMG data.
    \item Designed the pipeline for use by researchers outside computer science.
    \item Involved redesign of dataset acquisition and processing, \\
    network architecture setup, and augmentation handling.
    \item Pipeline based on PyTorch-Lightning.
    \item Using datasets like: NinaPro, SEED, and emg2pose.\\
\end{itemize}

\noindent \textbf{Vision Imitation Learning for Modular Robots} \textit{Fraunhofer Italia} \href{https://github.com/RobinU434/TaskParameterizedGaussianMixtureModels}{\faGithubSquare}\hfill May 2023 $-$ August 2023
\begin{itemize}[noitemsep,nolistsep,leftmargin=*]
    \item Internship project done under Dr. Andrea Giusti and \\ supervised by Dr. Marco Todescato.
    \item Implementation of a computer vision pipeline for \\ 
    3D hand trajectory tracking, object recognition, and \\
    segmentation with multi-stereo camera feed.
    \item Re-implementation of task-parameterized Gaussian mixture models based on \\ work from Sylvain Calinon (\href{https://calinon.ch/papers/Calinon-JIST2015.pdf}{paper}) for imitation learning.
    \item Development of Modular-Robot optimization pipeline.
    \item Web-based interface design for containerized pipeline setup.  
    \item Used technologies: Python, Numpy, OpenCV, ROS2, C++, Moveit, and Eigen.\\
\end{itemize}



\noindent \textbf{PACT} \textit{Deep Energy Retrofit} \hfill August 2022 $-$ March 2023
\begin{itemize}[noitemsep,nolistsep,leftmargin=*]
    \item Project done at Fraunhofer IPM under Dr. Annette Schmitt.
    \item Collected 3D data for point cloud segmentation with laser scanners. \\
\end{itemize}


\noindent \textbf{Amphora} \textit{Automated Pavement Operation} \hfill  September 2021 $-$ November 2022
\begin{itemize}[noitemsep,nolistsep,leftmargin=*]
    \item Project done with Fraunhofer-IPM team under Prof. Alexander Reiterer, \\
     with Dr. Annette Schmitt and international project-partners.
    \item Implemented multilabel image classification pipeline on laser scan data \\
    with convolutional neural networks (ResNet \& VGG) and Vision-Transformer.
    \item Conducted classification-oriented real-time and classification study with \\
    hyper-parameter-optimization (Bayesian-Optimization + Hyperband).
    % \item Conducted a study on real-time classification capabilities for image classification.
    \item Implemented image processing pipeline: laser scanner $\to$ geojson \\
\end{itemize} 


\noindent \textbf{Molecular Tumor Board Research Database}\hfill  March 2020 $-$ May 2021
\begin{itemize}[noitemsep,nolistsep,leftmargin=*]
    \item Project done with the IBSM team under Prof. Melanie Börries.
    \item Development of a database application oriented \\ 
    for molecular tumor research.
    \item Conducted a study on database design choices between \\
    SQL, document-based, and graph-based database systems.
    \item Knowledge acquisition of data collection from tissue, \\
    to sequencing to the final annotation pipeline. 
    \item Collecting wet-lab experience and testing UI-design\\
    choices together with the research staff.\\
\end{itemize}



\vspace{0.5cm}
\noindent \textbf{Private} \\

\noindent \textbf{LazySlurm} \textit{Software Project} \href{https://github.com/RobinU434/LazySlurm}{\faGithubSquare} \hfill since November 2025
\begin{itemize}[noitemsep,nolistsep,leftmargin=*]
    \item Terminal user interface for monitoring Slurm HPC jobs, \\
    published on PyPI as \texttt{lazyslurm-py} under the MIT license.
    \item Live overview of running and past jobs, log inspection, \\
    resource and CPU/GPU usage, job cancellation, editing and resubmission, \\
    and partition load.
    \item Remote mode running over a single background SSH connection \\
    with one-time two-factor authentication.
    \item Built in Python with the Textual framework.\\
\end{itemize}

\noindent \textbf{Simplified RNNs} \textit{Research Project} \hfill since April 2025
\begin{itemize}[noitemsep,nolistsep,leftmargin=*]
    \item Based on the GateL0rd paper written by Dr. Christian Gumbsch.
    \item Simplifying the architecture and conducting a study on latent state similarities.\\
\end{itemize}


\noindent \textbf{SPARK} \textit{Research Project} \href{https://github.com/RobinU434/SPARK}{\faGithubSquare} \hfill April 2025 $-$ November 2025
\begin{itemize}[noitemsep,nolistsep,leftmargin=*]
    \item Stochastic Policies Augmented for Robust Knowledge/Exploration.
    \item Builds on top of the idea of my Bachelor's thesis.
    \item Augmentation of stochastic policies for robust exploration and \\
    more expressive action distributions.
    \item Extension of SAC to leverage non-diagonal Gaussian distributions \\
    for action sampling.
    \item Benchmarking done on a 2D inverse kinematics environment.
    \item Using Stable-Baselines3 as the baseline for SAC. \\
\end{itemize}


\noindent \textbf{Tide $-$ Torch Integrated Data Engine} \textit{Software Project} \href{https://github.com/RobinU434/Tide}{\faGithubSquare}\hfill since February 2025
\begin{itemize}[noitemsep,nolistsep,leftmargin=*]
    \item Data toolkit for handling sequential data \\
    for single or multi-file datasets. 
    \item Compatible with PyTorch dataset, data-loader integration, \\
    feature engineering, time-based splitting, filtering, \ldots.
    \item Integration into publicly available time series data: \\
    EMG, climate, finance, robotics \ldots\\
\end{itemize} 


\noindent \textbf{Config2Class} \textit{Private} \href{https://github.com/RobinU434/Config2Class}{\faGithubSquare}\hfill since October 2024
\begin{itemize}[noitemsep,nolistsep,leftmargin=*]
    \item Developed Python tool to automatically generate structured configuration classes\\
    from various file formats.
    \item Adds support for code completion in IDE.
    \item Integrates with Hydra and OmegaConf.\\
\end{itemize}
 


\noindent \textbf{PyArgWriter (CLIfy)} \textit{Private} \href{https://github.com/RobinU434/PyArgWriter}{\faGithubSquare}\hfill since September 2023
\begin{itemize}[noitemsep,nolistsep,leftmargin=*]
    \item Developed a Python tool to automatically create command-line \\
    interface for Python projects based on the OOP paradigm.
    \item Implemented automatic code generation, code parsing tooling,\\ 
    tree-structured command generation, Hydra integration.\\
\end{itemize} 

\noindent \textbf{Hexapod Robot} \textit{Private} \hfill since May 2021
\begin{itemize}[noitemsep,nolistsep,leftmargin=*]
    \item Building custom Hexapod robot using the Freenove Big-Hexapod robot kit, \\
    Jetson Nano B01, Stereo vision, and custom stationary power supply. 
    \item Implementing RViz2 and Gazebo simulation plugins, ROS2 architecture, \\
    low-level servo controller, high-level walking controller. 
    \item Designing custom proprioceptive sensing for closed-loop \\
    controller design for cheap-PWM sensors.
    \item Conducted a study on mobile neural-depth estimation for stereo-vision. \\
\end{itemize}

%%%%%%%%%%%%%%%%%%%%%%%%%%%%%%%%%%%%%%%%%%%%%%%%%%%%%%%%%%%%%%%
% WORK EXPERIENCE
% What did you do? -> Project goals OR what problem did you solve?
% How did you do it? -> Skills and technologies
% What impact did you create? -> Numbers and percentages.
% Example: 
% + Developed an app for matching mentor and mentees for Android and iOS platform.
% + Successfully matched 85% of the applications and randomized the rest.
% 
% Talk about team work, initiative, soft skills.
%
% Can also include personal projects, competitions, contribution to Open source.
%%%%%%%%%%%%%%%%%%%%%%%%%%%%%%%%%%%%%%%%%%%%%%%%%%%%%%%%%%%%%%%
\noindent \textbf{\underline{ACADEMIC WORK EXPERIENCE}} \\
% \noindent \textbf{Metal Works Schick, metalwork} \hfill Aspach, Germany \\
% \textit{Internship} \hfill June 2018 $-$ July 2018
% \begin{itemize}[noitemsep,nolistsep,leftmargin=*]
%     \item {Collected experience in mechanical design}
%     \item {Collected experience in metal processing\\}
% \end{itemize}% 

% \noindent \textbf{Murrplastik GmbH, injection moulding} \hfill Sulzbach, Germany \\
% \textit{Internschip} \hfill July 2018
% \begin{itemize}[noitemsep,nolistsep,leftmargin=*]
%     \item {Collected experience in design of thermoplastic resin products}
%     \item {Gained knowledge about thermoplastic characteristics\\}
% \end{itemize}% 

% \noindent \textbf{Jagic Prototype Construction, } \hfill Mittelbrüden, Germany\\
% \textit{Internship} \hfill April 2019
% \begin{itemize}[noitemsep,nolistsep,leftmargin=*]
%     \item {Collected experience in  prototypes design}
%     \item {Assist in manufacturing process}
%     \item {Collected experience in detailed metal processing\\}
% \end{itemize}


\noindent \textbf{University of North Carolina at Chapel Hill}  \hfill Chapel Hill, NC, USA \\
\noindent Group of Prof. Weitong Zhang \hfill November 2026 $-$ December 2026\\
\textit{Visiting Researcher (upcoming)}
\begin{itemize}[noitemsep,nolistsep,leftmargin=*]
    \item Research stay on offline reinforcement learning \\
    with expressive probabilistic samplers.\\
\end{itemize}

\noindent \textbf{University of Tübingen}  \hfill Tübingen, Germany \\
\noindent Autonomous Learning Group  \hfill October 2025 $-$ March 2026\\
\textit{Student Assistant}
\begin{itemize}[noitemsep,nolistsep,leftmargin=*]
    \item Worked on humanoid robots in the group of Prof. Georg Martius, \\
    supervised by Jens Beißwenger and Anna Manasyan.
    \item Development of VLA models applied to humanoid robotics.\\
\end{itemize}

\noindent \textbf{BG Clinic - Tübingen}  \hfill Tübingen, Germany \\
\noindent Hand, plastic, reconstructive, and burn surgery  \hfill August 2024 $-$ February 2025\\
\textit{Student Assistant}
\begin{itemize}[noitemsep,nolistsep,leftmargin=*]
    \item Developed software interfaces in the group of \\
    Dr. Cosima Prahm and Dr. Dominik Steiner, \\
    supervised by Daniel Höglinger.
    \item Implemented a deep learning pipeline for EMG-to-pose data.
    \item Implemented convolutional, recurrent, and \\transformer architectures for time series data.
    \item Made the NinaPro dataset accessible through PyTorch.\\
\end{itemize}

\noindent \textbf{Volatilis}  \hfill Ludwigsburg, Germany \\
\noindent Software Development \hfill July 2024 $-$ February 2025\\
\textit{Freelancer}
\begin{itemize}[noitemsep,nolistsep,leftmargin=*]
    \item Implemented a browser-based simulation engine for risk modeling.
    \item Contributed to business and operational development. \\
\end{itemize}

\noindent \textbf{AI-Center Tübingen}  \hfill Tübingen, Germany \\
\noindent Central Office  \hfill December 2023 $-$ May 2024 \\
\textit{Student Assistant}
\begin{itemize}[noitemsep,nolistsep,leftmargin=*]
    \item Wrote a wiki covering organizational topics for \\
    international PhD students and post-doctoral researchers.
    \item Performed and analyzed a department-wide survey on \\ the experiences of international researchers.\\
\end{itemize}

\noindent \textbf{Fraunhofer Italia}  \hfill Bolzano, Italy \\
\noindent Group: Advanced Robotics  \hfill April 2023 $-$ August 2023 \\
\textit{Internship}
\begin{itemize}[noitemsep,nolistsep,leftmargin=*]
    \item Internship under the guidance of Dr. Andrea Giusti \\ and Dr. Marco Todescato.
    \item Built a vision-guided imitation learning pipeline for a \\ modular robot manipulator.
    \item Handled multi-camera streams for \\trajectory recording in assembly work.
    \item Designed interfaces for modular software components. \\
\end{itemize}

\noindent \textbf{Institute of Intelligent Machine-Brain Interfacing Technology (IMBIT)}  \hfill Freiburg im Breisgau, Germany \\
\noindent Group: Neurorobotics  \hfill July 2022 $-$ April 2023\\
\textit{Student Assistant}
\begin{itemize}[noitemsep,nolistsep,leftmargin=*]
    \item Worked in the group of Prof. Joschka Boedecker, \\
    supervised by Jasper Hoffmann.
    \item Maintenance and mechanical work for the Unitree A1 robot.
    \item Interface documentation for custom quadrupedal \\
    policy deployment on the Unitree A1.
    \item Literature review in reinforcement learning and control theory. \\
\end{itemize}

\noindent \textbf{Fraunhofer Institute for Physical Measurement Technology (IPM)}  \hfill Freiburg im Breisgau, Germany \\
\noindent Department: Object and Shape Detection  \hfill June 2021 $-$ March 2023\\
\textit{Student Assistant}
\begin{itemize}[noitemsep,nolistsep,leftmargin=*]
    \item Worked under Prof. Alexander Reiterer, supervised by \\
    Dr. Philipp von Olshausen and Dr. Annette Schmitt.
    \item Development of a multi-label image classification \\
    pipeline with convolutional neural networks.
    \item Conducted hyperparameter optimization study on \\ 
    fine-tuning of convolutional neural networks. 
    \item Collection of 3D point cloud data of urban habitation.\\
\end{itemize}

\noindent \textbf{University Clinic Freiburg im Breisgau}  \hfill Freiburg im Breisgau, Germany \\
\noindent Group: Institute of Medical Bioinformatics and System Biology (IBSM)  \hfill March 2020 $-$ May 2021 \\
\textit{Student Assistant}
\begin{itemize}[noitemsep,nolistsep,leftmargin=*]
    \item Developed a database application for molecular cancer \\
    research under Prof. Melanie Börries, supervised \\
    by Dr. Patrick Metzger and Maria Hess.
    \item Conducted a technical feasibility study involving \\
    the use of relational and non-relational databases.
    \item Implemented a web interface for an SQL database \\
    with Flask and SQLAlchemy. 
    \item Making the application feasible for server-deployment\\
    in a containerized fashion.\\
\end{itemize}


%%%%%%%%%%%%%%%%%%%%%%%%%%%%%%%%%%%%%%%%%%%%%%%%%%%%%%%%%%%%%%%
% SKILLS: Important and relevant to the job you are applying for
%%%%%%%%%%%%%%%%%%%%%%%%%%%%%%%%%%%%%%%%%%%%%%%%%%%%%%%%%%%%%%%

\noindent \textbf{\underline{TECHNICAL SKILLS}}
\begin{multicols}{2}
\begin{itemize}
    \item Python (NumPy, PyTorch (Lightning), Gymnasium, SciKit-Learn, Pandas, SB3, PettingZoo, \ldots)
    \item C, C++, C\#
    \item (NO-)SQL-Databases
    \item Docker
    \item ROS, ROS2
    \item Git
    \item SVN
    \item Mechanical CAD (Autodesk Inventor)
    \item Embedded systems (Arduino, ESP)
    \item SCRUM
    \item SLURM
\end{itemize}
\end{multicols}
\vspace{0.4cm}

%Programming in Python (one and a half years of experience), Prog (years of experience), Skill 3 (years of experience) \\

%%%%%%%%%%%%%%%%%%%%%%%%%%%%%%%%%%%%%%%%%%%%%%%%%%%%%%%%%%%%%%%
% Extra Curricular Activities, Leadership, etc 
%%%%%%%%%%%%%%%%%%%%%%%%%%%%%%%%%%%%%%%%%%%%%%%%%%%%%%%%%%%%%%%
\noindent \textbf{\underline{AWARDS}}
\begin{itemize}[noitemsep,nolistsep,leftmargin=*]
    \item \textbf{Field Robot Event} \textit{3rd place} \hfill August 2025
    \item \textbf{System Design Project} \textit{3rd place} \hfill February 2020 \\
\end{itemize}

\vspace{0.3cm}
\noindent \textbf{\underline{EXTRA CURRICULAR ACTIVITIES}} \\

\noindent \textbf{Private Statistics Tutor} \hfill September 2025 $-$ October 2025
\begin{itemize}[noitemsep,nolistsep,leftmargin=*]
    \item Tutoring undergraduate students in statistics. \\
\end{itemize}

\noindent \textbf{Organizational Staff at EWRL 18} \hfill September 2025
\begin{itemize}[noitemsep,nolistsep,leftmargin=*]
    \item Volunteer at the 18th European Workshop on Reinforcement Learning \\
    in Tübingen.
    \item Assisting at registration, discussions, demonstrations, \ldots \\
\end{itemize}

\noindent \textbf{Organizer Machine Learning Reading Club} \hfill February 2025 $-$ June 2025
\begin{itemize}[noitemsep,nolistsep,leftmargin=*]
    \item Took part in organizing a machine learning journal club \\
    for computer science master students at Politecnico di Milano, \\
    within the association ``Polimi Data Scientists''.
    \item Tasks: selecting papers, organizing speakers, communicating \\
    with fellow students.\\
\end{itemize}

\vspace{0.3cm}
\noindent \textbf{\underline{LANGUAGES}} \\

German (native), English (C1/C2)
\vspace{0.3cm} \\


% \noindent \textbf{Activity/ Role} \hfill Month, Year $-$ Month, Year
% \begin{itemize}[noitemsep,nolistsep,leftmargin=*]
% \item {What did you do, how did you do it and what did you achieve? \\}
% \end{itemize}

%%%%%%%%%%%%%%%%%%%%%%%%%%%%%%%%%%%%%%%%%%%%%%%%%%%%%%%%%%%%%%%
% Other Skills: you can add all your other skills here.
% Continue to keep only relevant skills
%%%%%%%%%%%%%%%%%%%%%%%%%%%%%%%%%%%%%%%%%%%%%%%%%%%%%%%%%%%%%%%

\noindent \textbf{\underline{PERSONAL INTERESTS}} \\

% \begin{multicols}{1}
Archery (bronze medal at European Championships 2018), Running, Road cycling, Photography
% \end{multicols}

\end{document}

